
\documentclass[11pt]{letter}
\usepackage[a4paper,margin=0.9in]{geometry}
\usepackage{fontspec}
\IfFontExistsTF{Arial}{\setmainfont{Arial}}{\setmainfont{TeX Gyre Heros}}
\usepackage{hyperref}
\hypersetup{colorlinks=true,urlcolor=blue}
\setlength{\emergencystretch}{3em}
\sloppy
\signature{Hongmin Li and Bian Bian}
\address{AlphaScience Lab\\\url{https://alphascience-lab.com}\\Correspondence: \href{mailto:lihongmin@alphascience-lab.com}{lihongmin@alphascience-lab.com}; \href{mailto:bian.bian@kitasato-u.ac.jp}{bian.bian@kitasato-u.ac.jp}}
\date{2026-05-27}
\begin{document}
\begin{letter}{Editors\\Communications Biology}
\opening{Dear Editors,}

We are pleased to submit our manuscript, ``Source-resolved toxicogenomics reveals compartment-conditioned BDE-47 responses in human placental evidence,'' for consideration as an Article in \emph{Communications Biology}.

This study addresses a common interpretive problem in toxicogenomics: curated chemical-gene signs can compress heterogeneous biological readouts, causing extracellular vesicle cargo, cellular proteome and cellular mRNA evidence to appear as directly comparable signed relations. Using BDE-47 as a focused placental toxicology case, we reconstruct the source context behind these apparent sign conflicts.

We show that the dominant CTD decrease-oriented evidence block for BDE-47 does not represent broad intracellular protein downregulation, but rather decreased representation in placenta-derived exosomal protein cargo. We then identify a compact compartment-conditioned contrast involving ILK, NFKB1 and SNX17: these proteins decrease in exosomal cargo while showing higher cytotrophoblast proteome abundance in published SWATH-MS data. A bounded DIA-NN raw-pilot analysis supports the same direction for this proteome contrast. Independent cytotrophoblast mRNA data further indicate that the exosomal cargo decrease block is not simply mirrored by uniform cellular mRNA depletion.

The work provides a source-resolved model for interpreting toxicogenomic sign conflicts as compartment- and readout-conditioned evidence, rather than as simple database artifacts or direct mechanistic conclusions. We explicitly bound all claims. In particular, the proteome contrast is interpreted as directional evidence, not as author-level differential abundance.

The manuscript uses only public secondary datasets and derived analyses. Data and code availability statements, an AI-assisted tools statement, author contributions and competing interests are included in the manuscript. All authors have approved this submission, and the manuscript is not under consideration elsewhere.

We believe this focused evidence-reconstruction study will be of interest to readers working in placental biology, toxicogenomics, extracellular vesicles and source-aware computational biology.

\closing{Sincerely,}
\end{letter}
\end{document}
